% Chapter 7, typeset from lectures/L07: the README is the opener and the closing review, and each
% part of Appendix A is a section, in the same order. Appendix B, the exercises, is the last
% section.

\chapter{Timers}
\label{c7:ch}

{\large\itshape A counter with a target count, drawn by hand before it is written.\par}

\begin{chapterbox}{In this chapter}
\begin{itemize}
  \item Timers: a counter that raises a flag once it reaches a target count.
  \item Turning clock ticks into real time, and the equality comparator that does it.
  \item The same timer in VHDL, made reusable by a \code{generic} rather than by rewiring.
  \item Composing a design out of pre-built modules instead of rewriting their logic.
  \item \code{walking_led} demonstrated on the DE0-CV: a single lit bit walking a row of LEDs.
\end{itemize}
\end{chapterbox}

This chapter builds a timer twice and then puts it to work. The first is a 4-bit timer in
CircuitVerse: Chapter~\ref{c6:ch}'s counter, plus an XNOR equality comparator and the clear that
makes it repeat, watched as \code{timeout} fires once every eleven edges, for exactly one edge. The
second is the same timer in VHDL, as a reusable module with its target count as a \code{generic}.
Then comes \code{walking_led}: the timer just written, plus Chapter~\ref{c4:ch}'s
\code{reset_sync} and \code{button_sync} unchanged, plus a loop to rotate the pattern, and onto the
board. Almost no new logic: that part is about composing what already exists.

Before the timer in \secref{c7:a2} runs, predict which count \code{timeout} fires on, and why it is
every eleven edges rather than every ten. The off-by-one is the whole reason the module's generic
is named \code{TICK_COUNT} and not \code{PERIOD}.

\code{walking_led} is the payoff of the last four chapters: three modules, two of them carried
unchanged from Chapter~\ref{c4:ch} and one written in this chapter, composed into something you can
watch from the back of the room. A timer is a counter plus one comparison, and this chapter assumes
the counter, so be comfortable with \secref{c6:a1} before you start. The synchronizers are not
re-explained here beyond a short recap. Running the exercises' testbenches, including the
multi-file \code{ghdl -a} line \code{blinker} and \code{walking_led} need, is covered in
\secref{c2:app:c}, and the Quartus project setup, pin assignment, compilation and programming used
for the board demonstration are in Appendix~\ref{app:quartus}.

% Sections 7.1 to 7.5, from lectures/L07/appendix/a_timers.md: one section per part, A.1 to A.5.

\section{Timers: A Counter with a Target}
\label{c7:a1}
A \textbf{timer} answers a different question than the counter from \secref{c6:a1}. A counter
answers ``what is the current count?''; a timer answers ``has a certain amount of time passed?''

It does that with three parts: an internal counter invisible from outside, a target value
\code{TICK_COUNT} fixed per instance, and a comparison that when true clears the counter and pulses
an output flag \code{timeout} for one clock cycle.

Because the system clock ticks at a known, fixed frequency, counting ticks is the same as measuring
elapsed time. On the DE0-CV that is 50,000,000 ticks per second:

\begin{console}
TICK_COUNT = seconds x 50,000,000 - 1
\end{console}

\noindent The \code{- 1} is there because the counter runs \code{0} through \code{TICK_COUNT}
inclusive, so a full period is \code{TICK_COUNT + 1} cycles, not \code{TICK_COUNT}.
Exercise~\ref{c7:ex:3} turns on exactly that, so it is worth reading off this formula rather than
the round number beside it:
\begin{itemize}
  \item \code{TICK_COUNT = 49,999,999}: timeout once per second.
  \item \code{TICK_COUNT = 24,999,999}: twice per second.
  \item \code{TICK_COUNT = 4,999,999}: ten times per second.
\end{itemize}
The rest of this chapter writes the round \code{50,000,000} and \code{5,000,000} where the exact
cycle count does not matter: one tick in fifty million is between 0.02 and 0.2\,ppm at these
values, far below the board oscillator's own tolerance.

This is the inverse of a blocking \code{_delay_ms()}. A blocking delay stops the whole program to
wait; a timer runs permanently in parallel with everything else and simply raises a flag, and
nothing in the design is ever blocked waiting for it.

\section{Building a Timer by Hand in CircuitVerse}
\label{c7:a2}
Before writing any VHDL, build a timer as a gate network, exactly as you did for the combinational
circuits of Chapters~\ref{c1:ch} and~\ref{c2:ch} and the flip-flops of Chapter~\ref{c3:ch}. A timer
is small enough to draw in full, and drawing it is what makes \secref{c7:a3}'s VHDL read as a
description of something you already understand rather than a new idiom to memorize.

Build it at a size you can watch: a \textbf{4-bit counter} with \code{TICK_COUNT = 10}, which is
\code{1010}. The real design counts to 50,000,000; nothing about the structure changes, only the
width.

\textbf{Start from the counter you already built} in \secref{c6:a1}: a 4-bit register \code{Q0} to
\code{Q3} sharing one clock, an adder and a constant \code{1} with the sum fed back to the \code{D}
inputs, and nothing else. If you did not keep it, it is three elements and about five minutes to
redraw. Confirm it still counts and wraps before adding anything, so that when the timer misbehaves
later you know the counter underneath is sound.

Everything here adds \textbf{two things on top of that}: a comparison saying ``the count has
arrived'', and a clear that makes it repeat. That is the entire difference between a counter and a
timer.

\textbf{Add the comparison.} You want \code{timeout} high exactly when the four bits spell
\code{1010}, and the clean way to build ``are these two numbers equal?'' is one \textbf{XNOR gate
per bit}. An XNOR outputs \code{1} when its inputs are the \emph{same}, which is \secref{c1:a1}'s
XNOR column read as an equality test. Wire four, each comparing a counter output against the
corresponding target bit (\code{Q3} against \code{1}, \code{Q2} against \code{0}, \code{Q1} against
\code{1}, \code{Q0} against \code{0}), driving the target side from a CircuitVerse constant rather
than a wire you might later mistake for a signal. Then AND the four outputs together; that single
AND gate is \code{timeout}.

Writing $\odot$ for XNOR, which is the one operator \secref{c1:a2} gave no symbol to:
\[
  \bool{timeout} = (\bool{Q3} \odot 1) \cdot (\bool{Q2} \odot 0) \cdot (\bool{Q1} \odot 1)
    \cdot (\bool{Q0} \odot 0)
\]
All four bits must match at once, which is what ``the counter has reached \code{1010}'' means. This
is a \textbf{general equality comparator}, worth recognizing as its own building block: the same
four XNORs plus an AND compare against any 4-bit value, and only the constants change.

\textbf{Add the clear}, so the timer restarts rather than running on: feed \code{timeout} back to
force the counter to \code{0000} on the next edge, gating each flip-flop's \code{D} input so
\code{timeout = 1} overrides the adder's output. Without it the counter keeps going to \code{1111}
and wraps on its own, so \code{timeout} fires every sixteen edges rather than every eleven. Both
are periodic; only one has the period you asked for.

\textbf{Then confirm, one clock edge at a time:}
\begin{itemize}
  \item The counter advances \code{0000} through \code{1010}, and \code{timeout} is low for every
    count except \code{1010}. Watch the four XNOR outputs as it counts and notice how often three
    are high at once. Three matching bits is not a match; the AND gate is what insists on all four.
  \item On the edge after \code{timeout} fires the counter is back at \code{0000} and
    \code{timeout} is low, so it is high for exactly one clock cycle, never two.
  \item The full cycle is therefore eleven edges, counts \code{0} through \code{10} inclusive,
    repeating indefinitely. That is \secref{c7:a1}'s \code{TICK_COUNT + 1} off-by-one, and this is
    the easiest place in the book to see it, because you can count the edges by eye.
\end{itemize}

\bookimage[htb]{0.6}{lectures/L07/appendix/images/timer_circuit.png}{The comparison logic:
four XNORs against the constants \code{1}, \code{0}, \code{1}, \code{0}, ANDed into
\code{timeout}.}{c7:fig:comparator}

\noindent The drawing in Figure~\ref{c7:fig:comparator} is cropped to the part this section adds.
It labels the counter's bits \code{counter[3]} down to \code{counter[0]} where the text above calls
them \code{Q3} to \code{Q0}, and the increment feeding them is drawn as gate-level carry logic
rather than the single adder block \secref{c6:a1} used, because that is what CircuitVerse gives
you when you build it a bit at a time. The four XNORs, their constants \code{1}, \code{0},
\code{1}, \code{0}, and the AND that combines them are the whole of what is new here.

Two things are worth noticing before the VHDL, because both are what the language buys you:
\begin{itemize}
  \item \textbf{The comparator is hard-wired to one target.} Changing \code{TICK_COUNT} from
    \code{10} to \code{12} means going back into the drawing and flipping two constants. It is a
    small edit, but it is an edit to the \emph{circuit}, and every instance needs its own copy. In
    \secref{c7:a3} the target becomes a generic.
  \item \textbf{The width is hard-wired too.} Counting to 50,000,000 needs 26 flip-flops and a
    26-bit comparator, 26 XNORs across 52 inputs: the same circuit, an unreasonable drawing. In
    \secref{c7:a3} it is the same six lines of VHDL either way.
\end{itemize}
That is the trade this book makes constantly: draw the circuit once, at a size you can see, to
know what the tool is building for you, then let the tool build it at a size you could not have
drawn.

\section{The \code{timer} Module in VHDL}
\label{c7:a3}
The same circuit as a reusable module. The course builds it live, and Exercise~\ref{c7:ex:4} asks
you to write it yourself afterwards, so it is printed here in full rather than shipped as a file:

\begin{vhdlcode}
entity timer is
    generic(TICK_COUNT: natural := 50_000_000);
    port(clock, reset_s2_n, enable: in std_logic;
         timeout                  : out std_logic);
end entity;

architecture behaviour of timer is
-- Internal tick counter.
signal counter: natural range 0 to TICK_COUNT;
begin
    process(clock, reset_s2_n) is
    begin
        if (reset_s2_n = '0') then
            timeout <= '0';
            counter <= 0;
        elsif (rising_edge(clock)) then
            timeout <= '0';
            if (enable = '1') then
                if (TICK_COUNT > counter) then
                    counter <= counter + 1;
                else
                    timeout <= '1';
                    counter <= 0;
                end if;
            end if;
        end if;
    end process;
end architecture;
\end{vhdlcode}

\begin{itemize}
  \item \code{TICK_COUNT} is a \textbf{generic}, not a constant baked into the architecture, so the
    same entity serves any frequency by instantiating it with a different \code{generic map}. The
    internal counter's range is written in terms of it, so the register is sized to the target
    rather than to whatever an unconstrained \code{natural} would give, the same point
    \secref{c6:a1} makes about \code{natural range 0 to 15}. A generic can size a signal precisely
    because it is fixed before synthesis runs.
  \item \code{reset_s2_n} is the already-\emph{synchronized} reset (\secref{c7:a4}), never the raw
    asynchronous \code{reset_n}. A timer never touches a raw asynchronous input directly.
  \item \code{enable} \textbf{freezes} the timer rather than resetting it. While it is \code{'0'},
    \code{counter} holds its value and \code{timeout} stays low, and a disabled timer resumes from
    where it left off rather than restarting. Chapter~\ref{c8:ch}'s \code{fsm_led} depends on this.
  \item \code{timeout} is a single-cycle pulse. The \code{else} branch that sets it is reached on
    only one edge: the edge at which \code{counter} is already sitting at \code{TICK_COUNT} and is
    cleared back to \code{0}. On the following edge the unconditional \code{timeout <= '0';} runs
    with nothing after it to override, so the pulse is exactly one cycle wide.
\end{itemize}
Compare this against \secref{c7:a2}'s gate network. The logic is the same, and \secref{c7:a2}
already named what the generic buys. One difference it did not: \secref{c7:a2}'s \code{timeout} is
a combinational decode, high during the cycle the count reads \code{1010}, while this one is
registered, high one cycle later, during the cycle \code{counter} reads \code{0}. Same period,
different phase, and different glitch behaviour, which is the registered-versus-combinational point
\secref{c8:a5} returns to.

\section{Composing a Full Timer Circuit}
\label{c7:a4}
A timer alone only pulses \code{timeout}; something has to react to that pulse, and a button has to
control the timer. Reusing Chapter~\ref{c4:ch}'s double-flop synchronizer, a complete
button-controlled, timer-driven LED circuit looks like Figure~\ref{c7:fig:floorplan}.

\bookimage{1.0}{lectures/L07/appendix/images/led_toggle_timer.png}{The full circuit at a glance:
synchronized button, timer, and toggled LED, with the detail readable in the two closer views of
Figures~\ref{c7:fig:sync} and~\ref{c7:fig:toggle}.}{c7:fig:floorplan}

\noindent Read that one as a floor plan rather than as wiring. It is a wide circuit, so at page
size its labels are small; what it is for is showing how few pieces there are and roughly where
each one sits. The two closer views below are where the signal names are legible, and between them
they cover everything in it.

Left to right:
\begin{itemize}
  \item \textbf{Synchronization and edge detection} (Figure~\ref{c7:fig:sync}). \code{button_n}
    passes through two flip-flops for metastability protection, then a third lets the circuit
    compare ``now'' against ``previous'' and detect a falling edge, producing a one-cycle
    \code{button_edge_s2} pulse.
  \item \textbf{The timer} (\secref{c7:a3}), enabled or disabled by \code{timer_enabled}, which is
    itself toggled by \code{button_edge_s2}.
  \item \textbf{Toggling the LED} (Figure~\ref{c7:fig:toggle}). The timer's \code{timeout} pulse
    flips the LED's state each time it fires, and the LED is forced off whenever the timer is
    disabled or the circuit is reset.
\end{itemize}

\bookimage{0.9}{lectures/L07/appendix/images/button_sync_and_flank_detection.png}{The synchronizer
and falling-edge detector: \code{button_n} through two flip-flops and a third for the previous
value, ANDed into \code{button_edge_s2}, which toggles \code{timer_enabled}.}{c7:fig:sync}

\bookimage{0.45}{lectures/L07/appendix/images/led_toggle.png}{The LED toggle logic, driven by the
timer's \code{timeout} pulse and forced off while the timer is disabled.}{c7:fig:toggle}

\noindent Two patterns matter here, independent of this circuit:
\begin{itemize}
  \item \textbf{Every asynchronous input is synchronized before it touches any other logic}, and
    every enable/toggle signal derived from a button is a synchronous signal from that point on.
    You reuse this unchanged in \secref{c7:a5} and again for Chapter~\ref{c8:ch}'s state machines.
  \item \textbf{One clock, and everything else is an enable.} \code{timeout} paces the LED by being
    a one-cycle \emph{enable} on a flip-flop that runs on the system clock, not by being wired into
    a clock input. It is tempting to route a counter or timer output into a clock port to ``slow
    things down''; do not. A design with one clock is a design whose timing the tool can analyze,
    and it is the single most common mistake a programmer makes writing a baud generator.
\end{itemize}

\section{FPGA Implementation: A Walking-LED Shift Register}
\label{c7:a5}
\code{walking_led} is the first design that composes three chapters' building blocks into one
hardware-demonstrable circuit: a single lit LED walking along a row, one position per timer tick,
started and stopped with a button. The course builds it live, and Exercise~\ref{c7:ex:6} asks you
to write it yourself afterwards, so it is printed here in full rather than shipped as a file.

It writes almost no new logic, reusing unchanged \code{reset_sync} and \code{button_sync}
(\secref{c4:a5} and \secref{c4:a6}) to synchronize and edge-detect the button, \code{timer}
(\secref{c7:a3}) to pace the rate, and the shift register from \S\S\ref{c6:a2}--\ref{c6:a4} for
the walking itself. That is what the last four chapters have been accumulating: a small set of
modules that compose.

You write all three yourself: \code{reset_sync} and \code{button_sync} in Chapter~4's
Exercises~\ref{c4:ex:8} and~\ref{c4:ex:9}, \code{timer} in Exercise~\ref{c7:ex:4} of this chapter.
Copy your own into the exercise directory before building this. Reusing a module means having one,
and this is the first design that asks you for the ones you already built.

\begin{vhdlcode}
entity walking_led is
    generic(LED_COUNT : natural := 8;
            TICK_COUNT: natural := 25_000_000);
    port(clock, reset_n, button_n: in std_logic;
         led                     : out std_logic_vector(LED_COUNT-1 downto 0));
end entity;

architecture behaviour of walking_led is
signal reset_s2_n, button_edge_s2: std_logic;
signal button_n_v, button_edge_s2_v: std_logic_vector(0 downto 0);
signal shift_enable, shift_tick  : std_logic;
signal shift_reg: std_logic_vector(LED_COUNT-1 downto 0);

begin
    led <= shift_reg;

    button_n_v(0)     <= button_n;
    button_edge_s2 <= button_edge_s2_v(0);

    reset_sync1: entity work.reset_sync
        port map(clock, reset_n, reset_s2_n);

    button_sync1: entity work.button_sync
        generic map(1)
        port map(clock, reset_s2_n, button_n_v, button_edge_s2_v);

    timer1: entity work.timer
        generic map(TICK_COUNT)
        port map(clock, reset_s2_n, shift_enable, shift_tick);

    SHIFT_ENABLE_PROCESS: process(clock, reset_s2_n) is
    begin
        if (reset_s2_n = '0') then
            shift_enable <= '0';
        elsif (rising_edge(clock)) then
            if (button_edge_s2 = '1') then
                shift_enable <= not shift_enable;
            end if;
        end if;
    end process;

    SHIFT_PROCESS: process(clock, reset_s2_n) is
    begin
        if (reset_s2_n = '0') then
            shift_reg <= (0 => '1', others => '0');
        elsif (rising_edge(clock)) then
            if (shift_tick = '1') then
                shift_reg <= shift_reg(LED_COUNT-2 downto 0) & shift_reg(LED_COUNT-1);
            end if;
        end if;
    end process;
end architecture;
\end{vhdlcode}

\noindent Design choices worth noting:
\begin{itemize}
  \item \code{button_n_v}/\code{button_edge_s2_v} bridge a width difference only:
    \code{walking_led}'s button signals are single bits while \code{button_sync}'s ports are
    \code{COUNT}-wide vectors. The two concurrent assignments wire the scalar to element \code{0}
    and back. They cost no hardware.
  \item This is a PISO register (\secref{c6:a4}) in its simplest form: parallel-loaded exactly
    once, on reset, with a single lit bit at position 0, and from then on only shifting. There is
    no runtime \code{load}, because this design never loads a new pattern.
  \item It is a \textbf{circular} shift register. The shift is \secref{c6:a2}'s expression
    unchanged, so the lit bit walks toward the MSB; the only difference from a plain PISO is what
    feeds bit 0, which is the bit falling off the top rather than a \code{serial_in} from
    elsewhere. A one-shot ``shift a pattern out'' register becomes an endlessly repeating pattern
    generator for free, just by recirculating its own output.
  \item \code{shift_enable} is \secref{c7:a4}'s toggle flip-flop reused verbatim: the button's
    synchronized edge pulse flips a single bit, and that bit becomes the timer's \code{enable}.
  \item \code{shift_tick} is the timer's \code{timeout} pulse, reused as the shift register's clock
    enable rather than driving an LED toggle. Same building block, different consumer, and exactly
    the ``one clock, everything else is an enable'' rule from \secref{c7:a4}.
  \item Both \code{LED_COUNT} and \code{TICK_COUNT} are generics, so
    \code{generic map(4, 25_000_000)} gives a 4-LED row and \code{generic map(8, 5_000_000)} a
    faster 100\,ms step, without touching the module.
\end{itemize}
On the board, as the course demonstrates it, the ports are assigned in the Pin Planner
(Appendix~\ref{app:quartus}): \code{clock} to the 50\,MHz oscillator, \code{reset_n} and
\code{button_n} to two push buttons, and \code{led} to a row of onboard LEDs. Press the button once
to start the bit walking, again to freeze it.

You can watch the same behaviour without a board by running the reference testbench, which
overrides both generics so a full lap takes a handful of clock cycles:

\begin{shell}
cd lectures/L07/exercises/walking_led
cp ../../../L04/exercises/reset_sync/reset_sync.vhd .   # the three you wrote yourself
cp ../../../L04/exercises/button_sync/button_sync.vhd .
cp ../timer/timer.vhd .
ghdl -a --std=93 reset_sync.vhd button_sync.vhd timer.vhd \
                 walking_led.vhd walking_led_tb.vhd
ghdl -e --std=93 walking_led_tb
ghdl -r --std=93 walking_led_tb --assert-level=error --stop-time=10ms
\end{shell}

% Section 7.6, from lectures/L07/README.md: the objectives, the questions to test yourself, and the
% next lecture.

\needspace{10\baselineskip}
\section{Review}
\label{c7:sec:review}

\subsection*{What you should be able to do now}
After this chapter, you should be able to:
\begin{itemize}
  \item Explain what separates a timer from a plain counter, and why the difference is one
    comparison.
  \item Build a timer as a gate network: an equality comparator from one XNOR per bit, and the
    clear that makes it repeat rather than wrap.
  \item Implement a timer in VHDL as a reusable module, converting a period in seconds into a
    \code{TICK_COUNT}, and say why \code{timeout} is a one-cycle pulse rather than a level.
  \item Compose a design out of pre-built modules rather than rewriting them, and say why a button
    controlling a timer still has to pass through a synchronizer first.
\end{itemize}

\needspace{6\baselineskip}
\subsection*{Questions to test yourself}
You should be able to explain:
\begin{enumerate}
  \item In what sense is a timer ``just a counter with a comparison''?
  \item Why must the internal counter be cleared back to \code{0} on reaching its target, rather
    than left to keep counting?
  \item Why is \code{timeout} a single-cycle pulse rather than a level that stays high?
  \item In the CircuitVerse version, changing the target count means rewiring the comparator. What
    replaces that rewiring in the VHDL version, and when is its value fixed?
  \item Why is it still necessary to run a button through a synchronizer before using its edge to
    enable a timer, even though the button has nothing to do with counting?
  \item \code{walking_led} writes almost no logic of its own. Name the three modules it composes
    and which chapter each came from.
\end{enumerate}

\needspace{6\baselineskip}
\subsection*{Looking ahead}
Chapter~\ref{c8:ch} is about state machines:
\begin{itemize}
  \item Finite state machines, the last building block in the book, designed by hand first: state
    diagram, state table, Karnaugh-derived logic, and a gate network in CircuitVerse, exactly the
    way this chapter's timer was drawn before it was written.
  \item Then in VHDL, where an enumerated type and a \code{case} statement replace the state table
    and the Karnaugh maps entirely.
  \item A timer is itself a state machine with two states and a single purpose.
    Chapter~\ref{c8:ch} generalizes that to circuits with many states and many transitions between
    them.
\end{itemize}

% Section 7.7, from lectures/L07/appendix/b_exercises.md.

\section{Exercises}
\label{c7:sec:exercises}

\begin{aside}
\textbf{How to work these.} Each exercise is tagged with the kinds of work it asks for:
\kinds{Paper} by hand, \kinds{CircuitVerse} in the browser, \kinds{VHDL} in a module of your own.

\textbf{How to check your work.} Every exercise below that asks you to write a VHDL module ships a
self-checking testbench under \file{lectures/L07/exercises} in the companion repository. Write your
module in its \file{exercises/<module>} directory, using the entity name and the \textbf{port
order} the exercise specifies, then run it with GHDL; \secref{c2:app:c} has the three commands.

From this chapter on, designs reuse modules you wrote earlier, and those are not shipped a second
time: each exercise says which of your own files to copy in and prints the \code{ghdl -a} line that
needs them.

No FPGA board is needed for any exercise. The Quartus synthesis and board-programming steps are
demonstrated in the course, and Appendix~\ref{app:quartus} shows them for a reader with a board;
your job is to get the VHDL right, and the testbench is how you confirm it. Short answers to the
parts you work on paper are in Appendix~\ref{app:answers}; reference solutions to the VHDL are
published in the companion repository, under \file{lectures/L07/solutions}. Try each exercise
yourself before looking at either.
\end{aside}

% ------------------------------------------------------------------------------------------
\exerciseset{Timers by Hand}
\label{c7:set:1}

\exercise{Retargeting the timer by hand}{Paper, CircuitVerse}
\label{c7:ex:1}
Build a second timer by hand in CircuitVerse, following \secref{c7:a2}, but with a target count of
\code{TICK_COUNT = 12} instead of the chapter's \code{10}.

\begin{parts}
  \item Retarget the comparator:
    \begin{itemize}
      \item Write \code{12} in 4-bit binary.
      \item The four XNOR gates stay exactly where they are. Say which of their constant inputs
        change, and to what.
      \item State, in one sentence, why an equality comparator built from XNORs needs no gate
        added, removed or rewired to compare against a different value.
    \end{itemize}
  \item Build it:
    \begin{itemize}
      \item Reuse your 4-bit counter and comparator from \secref{c7:a2}, changing only the
        constants.
      \item Set the clock period to \code{1000 ms} and confirm \code{timeout} fires once every
        thirteen clock edges, for exactly one edge each time.
    \end{itemize}
  \item Now remove the clear, so that \code{timeout} no longer forces the counter back to
    \code{0000}:
    \begin{itemize}
      \item Predict, before running it, how often \code{timeout} fires now.
      \item Run it and confirm.
      \item Explain the result in terms of the counter's own wraparound (\secref{c6:a1}).
    \end{itemize}
  \item Deliberately break the comparator: disconnect one XNOR's output from the AND gate and tie
    that AND input high instead.
    \begin{itemize}
      \item Predict which counts now assert \code{timeout}, before simulating.
      \item Confirm by stepping the counter through a full cycle.
      \item This is what ``three of the four bits match'' looks like, and it is the most common way
        a hand-drawn comparator goes wrong.
    \end{itemize}
\end{parts}
\note[Tip] counting \code{0} through \code{12} inclusive is thirteen edges, not twelve.
\Secref{c7:a1}'s note on \code{TICK_COUNT + 1} is the same off-by-one, and it is worth getting
straight here, where you can count the edges by eye, rather than later at 50\,MHz.

\exercise{Reading the timer module}{Paper}
\label{c7:ex:2}
Answer each of the following in one or two sentences, referring to the \code{timer} module in
\secref{c7:a3}:

\begin{parts}
  \item Why is \code{timeout} a single-cycle pulse, rather than a signal that stays high once the
    target count is reached?
  \item While a running timer's \code{enable} input is held low, what happens to its internal
    counter from one clock edge to the next: does it reset, pause, or keep counting?
  \item Why does the \code{timer} take \code{reset_s2_n} (the synchronized reset) rather than the
    raw \code{reset_n} as its reset input?
  \item The CircuitVerse timer in Exercise~\ref{c7:ex:1} needed rewiring to change its target
    count, while the VHDL module needs no edit at all. What replaces the rewiring, and at what point
    is its value fixed?
\end{parts}

% ------------------------------------------------------------------------------------------
\exerciseset{Timers in VHDL}
\label{c7:set:2}

\exercise{Tick counts at 50\,MHz}{Paper}
\label{c7:ex:3}
The DE0-CV board's system clock is \code{50 MHz}. For each target below, calculate the
\code{TICK_COUNT} generic you would pass to the \code{timer} module from \secref{c7:a3}:

\begin{parts}
  \item A \code{timeout} pulse once per second.
  \item A \code{timeout} pulse every \code{250} ms.
  \item A \code{timeout} pulse twenty times per second (\code{20 Hz}).
\end{parts}
\note[Tip] The relationship between a period and its tick count is:
\[
  \bool{TICK\_COUNT} = \bool{seconds} \times 50{,}000{,}000 - 1
\]
The \code{- 1} is the same off-by-one Exercise~\ref{c7:ex:1} turned on: the counter runs \code{0}
through \code{TICK_COUNT} inclusive, so a period is \code{TICK_COUNT + 1} clock cycles, not
\code{TICK_COUNT}.

\exercise{Writing the timer}{VHDL, Paper}
\label{c7:ex:4}
Write the \code{timer} module itself, as an entity named \code{timer}.

The course builds it live and \secref{c7:a3} prints it in full, so this is not a puzzle. Write it
anyway, without copying: every design from here to the end of the book instantiates this module,
including both capstones, and it is worth having built the thing you are about to reuse six times.

The entity has one generic:

\begin{booktable}{@{}p{5.2em}p{3.8em}p{5.4em}L@{}}
  \thead{Generic} & \thead{Type} & \thead{Default} & \thead{Description} \\
  \midrule
  \code{TICK_COUNT} & \code{natural} & \code{50_000_000} & Ticks to count before a timeout. One
    second at the DE0-CV's \code{50 MHz} clock, rounded: by Exercise~\ref{c7:ex:3}'s rule the
    exact value is \code{49_999_999}, and 20\,ns of error in a second is not worth the less
    readable constant. \\
\end{booktable}

\noindent and these ports:

\begin{booktable}{@{}p{5.6em}p{4.8em}p{4.6em}L@{}}
  \thead{Port} & \thead{Direction} & \thead{Type} & \thead{Description} \\
  \midrule
  \code{clock} & in & \code{std_logic} & System clock. \\
  \code{reset_s2_n} & in & \code{std_logic} & Active-low, already-synchronized reset.
    Asynchronous: it clears the counter and \code{timeout} with no clock edge involved. \\
  \code{enable} & in & \code{std_logic} & Timer enable. While low the counter \textbf{holds} its
    value, and resumes from there rather than restarting. \\
  \code{timeout} & out & \code{std_logic} & One-cycle pulse, asserted every
    \code{TICK_COUNT + 1} cycles while enabled. \\
\end{booktable}

\noindent The architecture must:
\begin{itemize}
  \item Declare the internal counter as \code{natural range 0 to TICK_COUNT}, so the register is
    sized by the generic rather than left unconstrained.
  \item Use one process, sensitive to \code{clock} and \code{reset_s2_n}, following the
    clocked-process template from \secref{c3:a6}.
  \item Clear both the counter and \code{timeout} on reset, outside the clocked branch.
  \item Drive \code{timeout} low unconditionally at the top of the clocked branch, and high only on
    the edge that clears the counter. That default-then-override is what makes the pulse exactly
    one cycle wide rather than two.
\end{itemize}
Then answer, in a sentence each:
\begin{itemize}
  \item \code{enable = '0'} must \textbf{pause} the counter, not clear it. Suppose you cleared it
    instead. Predict what the reference testbench reports, and which of the two properties above it
    is measuring when it does. Then run it and see. (Chapter~\ref{c8:ch}'s \code{fsm_led} depends
    on the pause, which is why the testbench bothers.)
  \item The counter counts \code{0} through \code{TICK_COUNT} inclusive, so a period is
    \code{TICK_COUNT + 1} cycles. Where in your code is that ``+ 1'' actually written? It is not a
    literal anywhere.
  \item \Secref{c7:a2}'s hand-drawn \code{timeout} is a combinational decode of the count; yours is
    registered. Both have the same period. Which cycle does each one fire on, and why does the
    difference matter for anything downstream that samples \code{timeout}?
\end{itemize}

\modulefigure{timer}

\selfcheck name your entity \code{timer}, with generic \code{TICK_COUNT} (\code{natural}, default
\code{50_000_000}), inputs \code{clock}, \code{reset_s2_n}, \code{enable}, and output
\code{timeout}, declared in that order; its testbench is in \file{exercises/timer}. It runs three
full periods, so a one-shot passes nothing, and pauses mid-count to confirm the count resumed rather
than restarted.

\textbf{Keep this file.} Exercise~\ref{c7:ex:5} below reuses it, so does \code{walking_led} in
Exercise~\ref{c7:ex:6}, and so do both Chapter~\ref{c8:ch} capstones. Every later exercise that
needs a timer asks you to copy this one in.

\exercise{A one-second blinker}{VHDL, Paper}
\label{c7:ex:5}
Write an entity named \code{blinker} that blinks an LED on and off once per second, by reusing the
\code{timer} module you just wrote rather than counting clock cycles itself.

The entity has one generic:

\begin{booktable}{@{}p{5.2em}p{3.8em}p{5.4em}L@{}}
  \thead{Generic} & \thead{Type} & \thead{Default} & \thead{Description} \\
  \midrule
  \code{TICK_COUNT} & \code{natural} & \code{25_000_000} & A half-second at the DE0-CV's
    \code{50 MHz} clock. Expose it as a generic rather than hard-coding the count, so a testbench
    can shorten the blink; the FPGA build just uses the default. \\
\end{booktable}

\noindent and these ports:

\begin{booktable}{@{}p{5.6em}p{4.8em}p{4.6em}L@{}}
  \thead{Port} & \thead{Direction} & \thead{Type} & \thead{Description} \\
  \midrule
  \code{clock} & in & \code{std_logic} & System clock. \\
  \code{reset_n} & in & \code{std_logic} & Active-low reset, raw from the outside world.
    Asynchronous: assert it and \code{led} goes out immediately, without waiting for a clock
    edge. \\
  \code{led} & out & \code{std_logic} & Cleared to \code{'0'} while the reset is asserted, and left
    there until the first timeout after the reset is released. \\
\end{booktable}

\noindent Note that this entity takes the \textbf{raw} \code{reset_n} rather than an
already-synchronized \code{reset_s2_n}. Synchronizing it is part of the exercise, and it is what
\secref{c4:a5} means by ``exactly one module in a design is the exception'': here that module is
\code{blinker} itself.

The architecture must:
\begin{itemize}
  \item Instantiate your \code{reset_sync} from Chapter~4's Exercise~\ref{c4:ex:8}, and use its
    \code{reset_s2_n} output as the reset for everything inside.
  \item Instantiate the \code{timer} module from \secref{c7:a3}:
    \begin{itemize}
      \item Pass your own \code{TICK_COUNT} generic straight through to the timer's
        \code{TICK_COUNT}.
      \item Hold its \code{enable} input high.
    \end{itemize}
  \item Toggle \code{led} in a synchronous process whenever the timer's \code{timeout} pulse fires.
\end{itemize}
Do not reimplement the counting logic, or the synchronizer, yourself.

Instead:
\begin{itemize}
  \item Reuse \code{timer} and \code{reset_sync} unchanged, via a \code{generic map} and
    \code{port map}, following the composition pattern in \secref{c7:a4}.
  \item Explain why a \textbf{half}-second timer period produces a \textbf{one}-second blink cycle
    (one full on-then-off period).
\end{itemize}
Then answer, in a sentence each:
\begin{itemize}
  \item The LED must be off during reset, and the testbench checks that it goes off with no clock
    edge in between. Which half of the assert-async-release-sync pattern does that check, and which
    half does it not?
  \item Your \code{timer} is reset by \code{reset_s2_n}, two edges after \code{reset_n} is
    released. What would go wrong if you fed it the raw \code{reset_n} instead? Name the failure,
    not just the rule.
\end{itemize}

\modulefigure{blinker}

\selfcheck name your entity \code{blinker}, with generic \code{TICK_COUNT} (\code{natural}, default
\code{25_000_000}), inputs \code{clock}, \code{reset_n}, and output \code{led}, declared in that
order; its testbench is in \file{exercises/blinker}. \textbf{Copy your own \file{reset_sync.vhd}
and \file{timer.vhd} into that directory} before you build; neither is shipped there, because you
wrote both. It overrides \code{TICK_COUNT} with a small value and checks that the LED toggles
repeatedly, holds for exactly one timer period between toggles, and that the first toggle arrives
late by your synchronizer's two edges (see \secref{c2:app:c}):

\begin{shell}
cd lectures/L07/exercises/blinker
cp ../../../L04/exercises/reset_sync/reset_sync.vhd .   # the two you wrote yourself
cp ../timer/timer.vhd .
ghdl -a --std=93 reset_sync.vhd timer.vhd blinker.vhd blinker_tb.vhd
ghdl -e --std=93 blinker_tb
ghdl -r --std=93 blinker_tb --assert-level=error --stop-time=10ms
\end{shell}

% ------------------------------------------------------------------------------------------
\exerciseset{The Walking LED}
\label{c7:set:3}

A single lit LED walking along a row, one position per timer tick, started and stopped with a
button. The course builds it live, and \secref{c7:a5} writes it up.

This is the payoff of the last four chapters, and it writes almost no logic of its own: three
modules, two of them carried unchanged from Chapter~\ref{c4:ch} and one written in this chapter,
composed. Exercise~\ref{c7:ex:6} builds it, and Exercises~\ref{c7:ex:7} and~\ref{c7:ex:8} modify
what you built.

\exercise{The walking LED}{VHDL, Paper}
\label{c7:ex:6}
Write an entity named \code{walking_led}.

Like the \code{timer} of Exercise~\ref{c7:ex:4}, the course builds it live and \secref{c7:a5}
prints it in full, so this is not a puzzle either. Write it anyway, without copying: composing
three modules you wrote yourself is the whole point of the exercise, and the capstones in
Chapter~\ref{c8:ch} ask for it again at twice the size.

The entity has two generics:

\begin{booktable}{@{}p{5.2em}p{3.8em}p{5.4em}L@{}}
  \thead{Generic} & \thead{Type} & \thead{Default} & \thead{Description} \\
  \midrule
  \code{LED_COUNT} & \code{natural} & \code{8} & How many LEDs the bit walks along. The row is this
    wide, and the walk wraps after this many steps. \\
  \code{TICK_COUNT} & \code{natural} & \code{25_000_000} & Passed straight through to your
    \code{timer}, so the bit moves one position per timer period: a half-second step at the
    DE0-CV's \code{50 MHz} clock. \\
\end{booktable}

\noindent and these ports:

\begin{booktable}{@{}p{4.6em}p{4.8em}p{10.8em}L@{}}
  \thead{Port} & \thead{Direction} & \thead{Type} & \thead{Description} \\
  \midrule
  \code{clock} & in & \code{std_logic} & System clock. \\
  \code{reset_n} & in & \code{std_logic} & Active-low reset, raw from the outside world.
    Synchronizing it is your job, as it was in \code{blinker}. \\
  \code{button_n} & in & \code{std_logic} & Active-low push button, raw. Each press starts or stops
    the walk. \\
  \code{led} & out & \code{std_logic_vector(}\allowbreak\code{LED_COUNT-1 downto 0)} & The LED
    row. Exactly one bit is lit at any time; on reset it is bit \code{0}. \\
\end{booktable}

\noindent The architecture composes three modules you have already written, and adds two
processes:
\begin{itemize}
  \item Instantiate your \textbf{\code{reset_sync}}, and use its \code{reset_s2_n} as the reset for
    everything else in the design, including the two instances below.
  \item Instantiate your \textbf{\code{button_sync}} with \code{generic map(1)}, for a one-cycle
    pulse per press. Its button ports are \code{COUNT}-bit vectors while \code{button_n} here is a
    scalar, so bridge the width with a one-element \code{std_logic_vector(0 downto 0)} in each
    direction. \Secref{c7:a5} shows this.
  \item Instantiate your \textbf{\code{timer}} with your own \code{TICK_COUNT}, driving its
    \code{enable} from the shift-enable flag below and taking its \code{timeout} as the shift tick.
  \item A process toggling \textbf{\code{shift_enable}} on each button pulse, cleared to \code{'0'}
    on reset. The walk therefore starts stopped.
  \item A process holding the \textbf{shift register} that drives \code{led}: on reset, load a
    single lit bit at position \code{0} and clear the rest; on each timer \code{timeout}, rotate one
    step toward the MSB, wrapping the top bit back around into bit \code{0}. That is a PISO
    register (\secref{c6:a4}) with its own serial output wired back to its serial input, which is
    what turns a one-shot shift-out into an endless walk.
\end{itemize}
Then answer, in a sentence each:
\begin{itemize}
  \item \code{shift_enable} drives the timer's \code{enable}, so stopping the walk stops the timer
    too. Exercise~\ref{c7:ex:4} made you show that \code{enable} \textbf{pauses} the count rather
    than clearing it. What does a reader see on the board, on the press that restarts the walk, if
    your timer clears instead?
  \item \code{button_sync} is instantiated with \code{COUNT = 1} here and with \code{COUNT = 2} in
    Chapter~\ref{c8:ch}'s \code{fsm_led}, from the same file with no edit. What would you have had
    to write instead if the width were fixed in the module rather than passed as a generic?
\end{itemize}

\modulefigure{walking_led}

\selfcheck name your entity \code{walking_led}, with generics \code{LED_COUNT} (\code{natural},
default \code{8}) and \code{TICK_COUNT} (\code{natural}, default \code{25_000_000}), inputs
\code{clock}, \code{reset_n}, \code{button_n}, and output \code{led}, declared in that order; its
testbench is in \file{exercises/walking_led}. It overrides the generics (\code{LED_COUNT = 4},
\code{TICK_COUNT = 3}) so a full lap takes a handful of clock cycles rather than two seconds, and it
checks the reset pattern, that nothing moves before the first press, and that the shift count over a
fixed window matches the pace the timer sets. Copy in the three modules it composes first:

\begin{shell}
cd lectures/L07/exercises/walking_led
cp ../../../L04/exercises/reset_sync/reset_sync.vhd .   # the three you wrote yourself
cp ../../../L04/exercises/button_sync/button_sync.vhd .
cp ../timer/timer.vhd .
ghdl -a --std=93 reset_sync.vhd button_sync.vhd timer.vhd \
                 walking_led.vhd walking_led_tb.vhd
ghdl -e --std=93 walking_led_tb
ghdl -r --std=93 walking_led_tb --assert-level=error --stop-time=10ms
\end{shell}

\exercise{Changing the walking LED}{VHDL, Paper}
\label{c7:ex:7}
Make the following changes one at a time, re-running the testbench after each. Parts \textbf{a)}
and \textbf{b)} change only the testbench's constants --- that they need no edit to your
\file{walking_led.vhd} at all is the point of them. Part \textbf{c)} is the one that touches the
module.

\begin{parts}
  \item Change the walking speed:
    \begin{itemize}
      \item Edit the \code{TICK_COUNT} constant in the \textbf{testbench} from \code{3} to
        \code{6}.
      \item It should still pass, and this time because the testbench adapts with you: its
        observation window is a fixed \code{PULSES * 6} cycles, and it derives the shift count it
        expects from \code{TICK_COUNT}, one shift per \code{TICK_COUNT + 1} cycles.
      \item Now find where that adaptation stops meaning anything. Keep raising \code{TICK_COUNT}
        past the window length and work out what the expected count becomes, and what the check is
        then actually proving. A test that passes for the wrong reason is worth being able to
        recognize.
      \item Then explain why changing the tick count cannot break the \emph{design}, only slow it
        down.
    \end{itemize}
  \item Change the width:
    \begin{itemize}
      \item Edit the testbench's \code{LED_COUNT} constant from \code{4} to \code{8}.
      \item It should still pass, with no change to your \file{walking_led.vhd} at all. Explain
        what the module would have needed to look like for this to require an edit.
    \end{itemize}
  \item Load two lit bits on opposite sides of the register on reset (e.g.\ \code{"10001000"} for
    \code{LED_COUNT = 8}), and let them walk in lockstep:
    \begin{itemize}
      \item Identify the single line in your \file{walking_led.vhd} you need to change.
      \item Before running anything, \textbf{predict} what the reference testbench will report.
      \item Then run it and see. Explain the result: is the design wrong, or is the testbench
        asserting something that is no longer the specification?
    \end{itemize}
\end{parts}
\note[Tip] Both \code{LED_COUNT} and \code{TICK_COUNT} are generics, so parts \textbf{a)} and
\textbf{b)} need no change to the module's body at all. That is the entire point of a generic: one
module, many sizes, no edits.

\exercise{Reversing the walk}{VHDL, Paper}
\label{c7:ex:8}
The current design only ever walks in one direction (toward the MSB, wrapping the top bit back
around to the bottom). Modify your \file{walking_led.vhd} so that a \textbf{second} button reverses
the walking direction.

\begin{parts}
  \item Add a second, synchronized button input:
    \begin{itemize}
      \item Route it through the same \code{button_sync} instance, passing \code{2} for its
        \code{COUNT} generic (\code{generic map(2)}, positionally, as everywhere else in this
        book).
    \end{itemize}
  \item Add a \code{direction} signal:
    \begin{itemize}
      \item Toggle it on the new button's edge, exactly the way \code{shift_enable} is toggled by
        the first button.
    \end{itemize}
  \item Make the rotation in \code{SHIFT_PROCESS} depend on \code{direction}:
    \begin{itemize}
      \item When \code{direction = '0'}, rotate toward the MSB (the book's convention from
        \secref{c6:a2}, and what you wrote in Exercise~\ref{c7:ex:6}):
\begin{vhdlcode}
shift_reg <= shift_reg(LED_COUNT-2 downto 0) & shift_reg(LED_COUNT-1);
\end{vhdlcode}
      \item When \code{direction = '1'}, rotate the other way, toward bit 0:
\begin{vhdlcode}
shift_reg <= shift_reg(0) & shift_reg(LED_COUNT-1 downto 1);
\end{vhdlcode}
    \end{itemize}
    This is the one place in the book where a shift register deliberately runs against the
    convention: reversing direction is the entire point of the exercise. Note that the
    second expression is exactly the mirror image \secref{c6:a2} warns you to watch for in other
    people's code.
  \item Your entity's ports have now changed: \code{button_n} is a two-bit vector rather than a
    single bit.
    \begin{itemize}
      \item Explain why \file{walking_led_tb.vhd} can no longer be used against your version, and
        what precisely GHDL will complain about when you try.
      \item This is the ordinary cost of changing an interface, and it is worth feeling once: a
        testbench is written against a specific entity, and widening a port breaks every
        instantiation of it, including the ones you did not write.
      \item Describe, in prose, the sequence of presses you would use to convince yourself the
        design works, and what you would expect the LEDs to do at each step. That sequence is
        exactly what a testbench for the new entity would have to drive.
    \end{itemize}
\end{parts}

